% © 2026 Ing. David Jaroš — CC BY-NC-ND 4.0
%
% This work is licensed under a Creative Commons Attribution-NonCommercial-NoDerivatives
% 4.0 International License (CC BY-NC-ND 4.0).
%
% License History: Earlier drafts (up to v0.3) were released under CC BY 4.0.
% From v0.4 onward, all material is released under CC BY-NC-ND 4.0 to protect
% the integrity of the theoretical work during ongoing academic development.
%
% See LICENSE.md for full license text.

\ifdefined\INCLUDEMODE
\else
  \documentclass[12pt,a4paper]{article}
  \usepackage[a4paper, margin=2.5cm]{geometry}
  \usepackage{amsmath, amssymb, amsthm}
  \usepackage{mathtools}
  \usepackage{hyperref}
  \usepackage{xcolor}
  \usepackage{mdframed}
  \usepackage{booktabs}

  \newtheorem{theorem}{Theorem}[section]
  \newtheorem{lemma}[theorem]{Lemma}
  \newtheorem{proposition}[theorem]{Proposition}
  \newtheorem{conjecture}[theorem]{Conjecture}
  \theoremstyle{definition}
  \newtheorem{definition}[theorem]{Definition}
  \theoremstyle{remark}
  \newtheorem{remark}[theorem]{Remark}

  \newmdenv[backgroundcolor=yellow!20, linecolor=orange, linewidth=1.5pt]{gapbox}
  \newmdenv[backgroundcolor=green!15, linecolor=green!60!black, linewidth=1.5pt]{resultbox}
  \newmdenv[backgroundcolor=red!10, linecolor=red!60, linewidth=1.5pt]{deadendbox}
  \newmdenv[backgroundcolor=blue!8, linecolor=blue!50, linewidth=1.5pt]{insightbox}

  \title{\textbf{Gap A --- Approach A3: Winding on $T^2(\tau)$}\\
         \large One-Loop Effective Action with Both $R_t$ and $R_\psi$ Nonzero}
  \author{David Jaro\v{s}}
  \date{2026-03-06}

  \begin{document}
  \maketitle

  \begin{abstract}
  We investigate whether the effective loop coefficient
  $B_{\mathrm{eff}}(R_t, R_\psi)$ computed on the two-dimensional torus
  $T^2(\tau)$ with both radii $R_t$ and $R_\psi$ nonzero can reproduce
  $B_{\mathrm{phenom}} \approx 46.3$.  The approach links Open Problem~A
  (the $B$ coefficient) to Open Problem~B (the correction factor
  $R \approx 1.114$), and tests whether the two-torus geometry with a
  natural ratio $R_t/R_\psi = \sqrt{137}$ or $\sqrt{139}$ provides the
  required enhancement over $B_0 = 8\pi$.
  \end{abstract}
\fi

%──────────────────────────────────────────────────────────────────────────────
\section{Setup: Two-Dimensional Torus}
\label{sec:setup}
%──────────────────────────────────────────────────────────────────────────────

Complex time $\tau = t + i\psi$ lives on a two-torus $T^2$ with radii
$R_t$ (real-time circle) and $R_\psi$ ($\psi$-circle).  The current UBT
derivation of $B_0 = 8\pi$ sets $R_t \to \infty$ (non-compact real time)
and keeps only the $\psi$-circle.

When both circles are compact:
\begin{equation}
  \tau \;\sim\; \tau + 2\pi R_t, \qquad
  \tau \;\sim\; \tau + 2\pi i R_\psi.
\end{equation}
The modular parameter of $T^2$ is $\rho = i\,R_\psi/R_t$, which lies in
the upper half-plane.

The kinetic action on $T^2$ is:
\begin{equation}
  S_{\mathrm{kin}}[\Theta]
  \;=\; \frac{1}{2}\int_{T^2} \Bigl[
      \frac{1}{R_t^2}|\partial_t\Theta|^2
    + \frac{1}{R_\psi^2}|\partial_\psi\Theta|^2
  \Bigr]\,dt\,d\psi.
\end{equation}

%──────────────────────────────────────────────────────────────────────────────
\section{Winding Mode Spectrum on $T^2$}
\label{sec:spectrum}
%──────────────────────────────────────────────────────────────────────────────

The Fourier decomposition on $T^2$:
\begin{equation}
  \Theta(t,\psi)
  \;=\; \sum_{m,n\in\mathbb{Z}} \Theta_{m,n}\,
        e^{i m t / R_t}\,e^{i n \psi / R_\psi}.
\end{equation}
The mode $(m,n)$ has mass squared
\begin{equation}
  M_{m,n}^2 \;=\; \frac{m^2}{R_t^2} + \frac{n^2}{R_\psi^2}.
\end{equation}

The one-loop vacuum polarisation is (schematically):
\begin{equation}
  \Pi(q^2)
  \;\sim\; \sum_{m,n} \int \frac{d^4k}{(2\pi)^4}\,
           \frac{1}{(k^2 + M_{m,n}^2)^2}
\end{equation}
in Euclidean space.  The sum over the KK tower $(m,n)$ gives an effective
$B$ coefficient that depends on $R_t$ and $R_\psi$.

%──────────────────────────────────────────────────────────────────────────────
\section{Effective $B$ Coefficient on $T^2$}
\label{sec:effective_B}
%──────────────────────────────────────────────────────────────────────────────

Retaining only the zero-mode in real time ($m = 0$, same as flat-space
limit in the $t$ direction) but summing over all $\psi$-modes:
\begin{equation}
  B_{\mathrm{eff}}(R_\psi) \;=\; B_0 \;\times\; f(R_\psi M_{\mathrm{UV}}),
\end{equation}
where $f$ is a form factor from the KK sum (shown in
\texttt{tools/compute\_B\_KK\_sum.py} to fail for a single circle).

When \emph{both} circles are compact, new contributions arise from modes
with $m \neq 0$.  Under Poisson summation (see
\texttt{ubt\_with\_chronofactor/scripts/spectral/poisson\_duality\_demo.py}):
\begin{equation}
  \sum_{m,n\in\mathbb{Z}} e^{-\pi\alpha(m^2 R_t^{-2} + n^2 R_\psi^{-2})}
  \;=\; \frac{R_t R_\psi}{\alpha}
    \sum_{m,n\in\mathbb{Z}} e^{-\pi(m^2 R_t^2 + n^2 R_\psi^2)/\alpha}.
\end{equation}
This is the Poisson duality $W(\tau) = \tau^{d/2} S(\tau)$ for $d=2$.

The one-loop effective action on $T^2$ receives a \emph{Poisson-dual}
contribution proportional to $R_t R_\psi$.  Normalising by the flat-space
result gives:
\begin{equation}
  \frac{B_{\mathrm{eff}}(R_t, R_\psi)}{B_0}
  \;\approx\; 1 + \frac{R_t R_\psi}{\lambda_{\mathrm{UV}}^2}
              \times (\text{modular correction}),
\end{equation}
where $\lambda_{\mathrm{UV}}$ is the UV cutoff.

%──────────────────────────────────────────────────────────────────────────────
\section{Natural Ratios from Prime Geometry}
\label{sec:natural_ratios}
%──────────────────────────────────────────────────────────────────────────────

The UBT prime attractor theorem selects $p \in \{137, 139\}$.  Two natural
radius ratios arise:
\begin{equation}
  \frac{R_t}{R_\psi} \;=\; \sqrt{137}
  \quad\text{or}\quad
  \frac{R_t}{R_\psi} \;=\; \sqrt{139}.
\end{equation}

Under Poisson duality the winding-to-momentum mode ratio is:
\begin{equation}
  \frac{R_t}{R_\psi} \;=\; \sqrt{p^*}
  \quad\Longrightarrow\quad
  \rho \;=\; \frac{R_\psi}{R_t} \;=\; \frac{1}{\sqrt{p^*}}.
\end{equation}

The modular $j$-function at $\rho = i/\sqrt{p^*}$:
\begin{align}
  j(i/\sqrt{137}) &\approx -744 + 196884\,e^{-2\pi/\sqrt{137}} + \cdots\\
                  &\approx -744 + 196884 \times e^{-1.695} \approx -744 + 35540 \approx 34796.
\end{align}
This large number is not directly related to $B_{\mathrm{phenom}}$.

\begin{insightbox}
\textbf{Key computation:}
If $R_t/R_\psi = \sqrt{p^*}$ is the natural ratio, the Poisson-dual
correction to $B_0$ at one loop is:
\begin{equation}
  \delta B \;=\; B_0 \times 2\,K_0\!\Bigl(2\pi\sqrt{p^*}\Bigr)\,\sqrt{p^*},
\end{equation}
where $K_0$ is the modified Bessel function of the second kind.
For $p^* = 137$:
$K_0(2\pi\sqrt{137}) \approx K_0(73.5) \approx 10^{-32}$ ---
exponentially suppressed.

This confirms that the Poisson-dual winding contribution is negligible
when $R_t/R_\psi = \sqrt{p^*} \gg 1$.
\end{insightbox}

\begin{deadendbox}
\textbf{Dead end (large $R_t/R_\psi$ ratio):}
For $R_t/R_\psi = \sqrt{137}$, the Poisson-dual KK contribution to
$\delta B$ is exponentially suppressed ($\sim e^{-73}$).
This does not produce the required $B_{\mathrm{eff}} = 46.3$.
\end{deadendbox}

%──────────────────────────────────────────────────────────────────────────────
\section{Small-Ratio Case: $R_t = R_\psi$}
\label{sec:small_ratio}
%──────────────────────────────────────────────────────────────────────────────

For the self-dual torus $R_t = R_\psi = R$, the two-loop effective action
is invariant under $T$-duality $R \to \alpha'/R$.  The effective $B$
coefficient at a self-dual radius:
\begin{equation}
  B_{\mathrm{eff}}(R = 1) \;=\; B_0 \times Z_T,
\end{equation}
where $Z_T$ is the torus partition function.  For a single real boson on a
self-dual circle:
\begin{equation}
  Z_T(R=1) \;=\; \sum_{m,n} q^{(m+n)^2/4}\,\bar{q}^{(m-n)^2/4}
\end{equation}
at $q = e^{2\pi i \tau}$ with $\tau = i$ (self-dual point).  This sums to
$Z_T(R=1) = 1$ in the appropriate normalisation (standard string theory
result).

Thus $B_{\mathrm{eff}}(R_t = R_\psi) = B_0$ --- no enhancement.

%──────────────────────────────────────────────────────────────────────────────
\section{Connection to Open Problem B: $R \approx 1.114$}
\label{sec:openproblemB}
%──────────────────────────────────────────────────────────────────────────────

Open Problem B is the correction factor $R \approx 1.114$ in the UBT
fine-structure derivation.  If $B_{\mathrm{phenom}} = B_0 \times R$, then:
\begin{equation}
  R \;=\; \frac{B_{\mathrm{phenom}}}{B_0}
    \;=\; \frac{46.3}{8\pi}
    \;\approx\; 1.842.
\end{equation}
But $R \approx 1.114$ is a \emph{different} correction (documented in
\texttt{docs/PROOFKIT\_ALPHA.md~§5}).

Approach~A3 cannot close Gap~A via the torus geometry alone; it would
require $B_{\mathrm{eff}}/B_0 \approx 1.842$, not 1.114.  The two
correction factors appear to be independent.

%──────────────────────────────────────────────────────────────────────────────
\section{Conclusion}
\label{sec:conclusion}
%──────────────────────────────────────────────────────────────────────────────

\begin{deadendbox}
\textbf{Dead end (Approach A3, torus winding):}
The one-loop effective action for $\Theta$ on $T^2(R_t, R_\psi)$ does not
produce $B_{\mathrm{eff}} = 46.3$ for any natural ratio $R_t/R_\psi$.

\begin{itemize}
  \item $R_t/R_\psi = \sqrt{p^*}$: Poisson-dual correction $\sim e^{-73}$,
        negligible.
  \item Self-dual $R_t = R_\psi$: $B_{\mathrm{eff}} = B_0 = 8\pi$.
  \item $R_t \to \infty$ (current UBT): $B_{\mathrm{eff}} = B_0 = 8\pi$.
\end{itemize}

\textbf{Conclusion:} The winding correction on $T^2$ cannot bridge the
gap from $B_0 = 8\pi$ to $B_{\mathrm{phenom}} = 46.3$.
Approach~A3 is a \textbf{complete dead end}.

\textbf{Honest statement:} We have tested all three natural radius ratios.
None reproduces $B_{\mathrm{phenom}}$.  The torus winding approach is
definitively ruled out at one loop.
\end{deadendbox}

\ifdefined\INCLUDEMODE
\else
  \end{document}
\fi
